Learned reconstruction through scattering media reports steadily improving fidelity, yet a reconstruction can spectrally invent content the measurement cannot physically contain. We isolate a free preprocessing choice---the bandwidth of the training target---as a causal control of this fabrication. In a controlled speckle-imaging protocol with a compute-matched control (identical initialization and identical additional optimization), fine-tuning against bandwidth-matched targets reduces the reconstruction's out-of-band spectral energy to $0.34\pm0.11$ (matched) and $0.20\pm0.08$ (half) of the control across two architectures and three seeds (6/6 units, one-sided sign test $p\approx0.016$), while the no-reference forward-consistency endpoint changes by only $+0.026$ on a $0.36$--$0.40$ baseline. Raw fabrication and suppression strength both increase in the larger of the two tested capacities. A Wiener linear baseline constrains the fabrication budget ($2$--$6\times$ below the raw-target and matched arms) but stays under-fitted: with optimal scalar rescaling its relative forward error remains about $40\%$ above the learned arms. A cross-diffuser distance-law analysis finds no class boundary, supporting protocol-level rather than label-level reporting. A real-data assessment on a public experimental dataset quantifies system parameters but shows that quantitative reconstruction requires a dedicated forward-calibration stage. Throughout, out-of-band energy serves as a fabrication proxy, not a perceptual metric.
